\icode{E-MK-SI-L2-260928-001}
一个朝着某方向做直线运动的物体，在时间\(t\)内的平均速度是\(v\)，紧接着在时间\(\frac{t}{2}\)内的平均速度是\(\frac{v}{2}\)，则物体在这段时间内的平均速度是\paren[D]
\begin{choices}
	\item \(v\)
	\item \(\dfrac{2}{3}v\)
	\item \(\dfrac{4}{3}v\)
	\item \(\dfrac{5}{6}v\)
\end{choices}
\begin{answer}
	D\qquad \kaishu 由平均速度的定义可知\(\overline{v} = \frac{vt + \frac{v}{2}\cdot\frac{t}{2}}{t + \frac{t}{2}} = \frac{5}{6}v\)
\end{answer}